%%% algorithm

This study aims to accelerate a process to identify moving object candidates from a sequence of images. Many sophisticated applications are available to extract sources from images \citep[e.g., \texttt{DAOFIND} and \texttt{Source Extractor}; ][]{stetson_daophot_1987,bertin_sextractor:_1996}. The list of the extracted sources may contain fixed stars, cosmic rays, and artificial noises, as well as possible moving objects. It is not an easy task to isolate moving object candidates from such distractors. This process can be a bottleneck in moving object extraction.

We developed a dedicated algorithm, Tracklet Extraction Engine (\texttt{tracee}), which efficiently extracts moving object candidates from a source list. The \texttt{tracee} is originally developed for an asteroid search but is possibly utilized in other fields.


\subsection{Problem Setting}
The algorithm we develop, \texttt{tracee},  is intended to deal with short video data. Thus, we assume that the motions of moving objects are linear and uniform. Thus, the problem is described as extracting linearly aligned points from a point cloud in a three-dimensional space. A single data set may contain multiple moving objects. The algorithm is required to detect them separately. The moving objects possibly cross fixed stars and each other. The algorithm should work properly in such cases.

\textit{Hough transformation} \citep{hough_method_1962} a plausible algorithm to deal with the problem, where each source votes for points in the hyperspace where each point is associated with a line segment. Such voting-based algorithms are, however, usually memory-consuming, while some efficient implementations are proposed \citep{li_fast_1986}. Thus, we do not adopt Hough transformation. Instead, we split the problem into two parts; First, we construct a $k$-nearest neighbor graph of the sources and then extract linearly aligned edges from the graph. An overview of the procedures is presented in Algorithm~\ref{algorithm:procedures}. Detailed descriptions are described below.


\begin{algorithm}[tp]
  \caption{Overview of \texttt{tracee}}
  \label{algorithm:procedures}
  \DontPrintSemicolon
  \SetKwProg{Part}{}{}{}
  \Part{\textbf{\upshape Part 1:} $k$-NN Graph}{
    \Input{$V=\left\{v_i\middle|_{i=1{\ldots}N}\right\}$}
    \Output{$E=\left\{(v_i,v_j)\middle|_{i,j \in \{1{\ldots}N\}}\right\}$}
    {Construct a $k$-NN graph $G(V,\,E)$ from $V$.}\;
    {Remove inappropriate edges in $E$.}\;
    {Redirect the edges in the chronological order.}\;
    \KwRet{$E$}\;
  }
  \BlankLine
  \Part{\textbf{\upshape Part 2:} Line Segment Grouping}{
    \Input{$E=\left\{(v_i,v_j)\middle|_{i,j \in \{1{\ldots}N\}}\right\}$}
    \Output{$B=\left\{b_m\middle|_{m=1{\ldots}M}\right\}$}
    {Find linearly aligned edges in $E$.}\;
    {Construct a set of baselines $B$.}\;
    {Remove inappropriate baselines from $B$.}\;
    \KwRet{$B$}\;
  }
\end{algorithm}


\subsection{Procedures}
\subsubsection{Input and Output}
The \texttt{tracee} receives a list containing the source positions and the timestamps of extraction. Hereafter, we call an element of the list a \textit{vertex}. The input data set $V$ is described as follows:
\begin{equation}
  \label{eq:algo:sources}
  V = \Bigl\{
    v_i = \left\{ x_i, \, y_i, \, t_i \right\} \Bigm|
    i = 0,\ldots,N
    \Bigr\},
\end{equation}
where $i$ is the index of the source, $v_i$ is the vertex corresponding to the $i$th source, $x_i$ and $y_i$ are the source positions, $t_i$ is the extraction timestamp, and $N$ is the total number of the extracted sources.

The tracklet is defined as a line segment or an edge. We define a \textit{baseline} as a data structure containing a line segment and associated vertices:
\begin{equation}
  \label{eq:algo:tracklet}
    b = \left\{ v^s, v^e, w \right\},
\end{equation}
where $v^s$ and $v^e$ are respectively the starting and end points of the line segment, $w$ contains the list of the vertices associated with the baseline. The definitions of these elements are described below:
\begin{equation}
  \label{eq:algo:tracklet:detail}
  \begin{cases}
    v_s = \left\{ x^s,\, y^s,\, t^s \right\}, \\
    v_e = \left\{ x^e,\, y^e,\, t^e \right\}, \\
    w = \left\{ \theta_1,\, \theta_2,\, \ldots,\, \theta_{K} \right\},
  \end{cases}
\end{equation}
where $K$ is the number of the vertices associated with the baseline and $\theta_i$ denotes the index of the vertex. The algorithm is expected to return a set of baselines, $B$:
\begin{equation}
  \label{eq:algo:sources}
  B = \Bigl\{
    b_m = \left\{ v^\mathrm{s}_m,\, v^\mathrm{e}_m,\, w_m \right\}
    \Bigm| m = 0,\ldots,M
  \Bigr\},
\end{equation}
where $m$ is the index of the baseline and $M$ is the total number of the baselines.

\subsubsection{$k$-Nearest-Neighbor Graph}

\subsubsection{Line Segment Grouping}
\lipsum[5]


%%% Local Variables:
%%% mode: latex
%%% TeX-master: "../manuscript"
%%% End:
